% !TEX TS-program = xelatex
% Command for running this example (needs latexmkrc file):
%    latexmk -bibtex -pdf main.tex

%	نمونه پایان‌نامه آماده شده با استفاده از کلاس tehran-thesis، نگارش 1
%	سینا ممکن، دانشگاه تهران 
%	https://github.com/sinamomken/tehran-thesis
%	گروه پارسی‌لاتک
%	http://www.parsilatex.com
%	این نسخه، بر اساس نسخه‌ 0.1 از کلاس IUST-Thesis آقای محمود امین‌طوسی آماده شده است.
%        http://profsite.sttu.ac.ir/mamintoosi

%----------------------------------------------------------------------------------------------
% اگر قصد نوشتن پروژه کارشناسی را دارید، در خط زیر به جای msc، کلمه bsc و اگر قصد نوشتن رساله دکترا را دارید، کلمه phd را قرار دهید. کلیه تنظیمات لازم، به طور خودکار، اعمال می‌شود.

% اگر مایلید پایان‌نامه شما دورو باشد به جای oneside در خط زیر از twoside استفاده کنید.

% برای حاشیه‌نویسی و کم کردن صفحات ابتدایی، گزینه draft را وارد و برای نسخه نهایی آن را حذف کنید.

% برای استفاده از قلم‌های سری IR Series گزینه irfonts را وارد و برای استفاده از قلم‌های X Series 2 آن را حذف کنید.

\documentclass[
twoside
% ,openany
,phd
,irfonts
% ,draft
]{./tex/tehran-thesis}

% فایل commands.tex را مطالعه کنید؛ چون دستورات مربوط به فراخوانی بسته‌ها، فونت و دستورات خاص در این فایل قرار دارد.
\input{./tex/commands}

% مشخصات پایان‌نامه را در فایلهای faTitle و enTitle وارد نمایید.
% !TeX root=../main.tex
% مشخصات فارسی رساله دکتری
%%%%%%%%%%%%%%%%%%%%%%%%%%%%%%%%%%%%
\university{دانشگاه تهران}
\college{دانشکدگان علوم و فناوری‌های میان‌رشته‌ای}
\faculty{دانشکدگان علوم و فناوری‌های میان‌رشته‌ای}
% گروه آموزشی خود را وارد کنید (در صورت نیاز)
\department{گروه علوم داده}
% رشته تحصیلی خود را وارد کنید
\subject{مهندسی شبکه‌های کامپیوتری}
% گرایش خود را وارد کنید
\field{؟؟}
% عنوان رساله
\title{بازسازی کور مشترک کدهای توربو با بازیابی همزمان کدگذارهای RSC و درهم‌ریزهای داخلی و خارجی}
% عنوان کوتاه برای سربرگ صفحات (تا با عنوان فصل تداخل نکند)
\runningtitle{بازسازی کور مشترک کدهای توربو}
% نام استاد(ان) راهنما
\firstsupervisor{دکتر مهدی تیموری}
\firstsupervisorrank{دانشیار}
%\secondsupervisor{دکتر راهنمای دوم}
%\secondsupervisorrank{استادیار}
% نام استاد(ان) مشاور؛ در صورت نداشتن مشاور، خطوط زیر را غیرفعال کنید.
\firstadvisor{دکتر کاکائی}
\firstadvisorrank{استادیار}
%\secondadvisor{دکتر مشاور دوم}
% نام داوران داخلی و خارجی
\internaljudge{دکتر داور داخلی}
\internaljudgerank{دانشیار}
\externaljudge{دکتر داور خارجی}
\externaljudgerank{دانشیار}
\externaljudgeuniversity{دانشگاه داور خارجی}
% نماینده کمیته تحصیلات تکمیلی
\graduatedeputy{دکتر نماینده}
\graduatedeputyrank{دانشیار}
% نام دانشجو
\name{شهریار}
\surname{توقی}
\studentID{۰۰۰۰۰۰۰۰۰}
\thesisdate{شهریور ۱۴۰۵}

%%%%%%%%%%%%%%%%%%%%%%%%%%%%%%%%%%%%%%%%%%%%%%%%%%%%
\dedication
{
{\Large تقدیم به:}\\
\begin{flushleft}{
	\huge
	\hspace{1em}\\
}
\end{flushleft}
}

\acknowledgement{
سپاس خداوندگار حکیم را که با لطف بی‌کران خود، آدمی را به زیور عقل آراست.

در آغاز وظیفهٔ خود می‌دانم از زحمات بی‌دریغ استاد(ان) راهنما و مشاور صمیمانه تشکر و قدردانی کنم که در طول انجام این رساله همواره راهنما و مشوق اینجانب بودند.

از خانواده و دوستانی که در این مسیر همراهی کردند نیز سپاسگزارم.
}

%%%%%%%%%%%%%%%%%%%%%%%%%%%%%%%%%%%%
\fa-abstract{
شناسایی و بازسازی کور ساختار کدهای کانال، به‌ویژه کدهای توربو، در مخابرات غیرهمکارانه یک مسئلهٔ کلیدی است. در این رساله، مسئلهٔ بازسازی مشترک پلی‌نوم‌های کدگذارهای RSC، درهم‌ریز داخلی و درهم‌ریز خارجی با درنظرگرفتن خاتمهٔ ترلیس و نویز صورت‌بندی و حل می‌شود. رویکرد پیشنهادی بر قیود جبری کدگذار، آزمون آماری قیود تحت نویز و جست‌وجوی درختی مقید استوار است تا پارامترهای مجهول از روی مشاهدات سخت (و در تعمیم‌ها، نرم) بازیابی شوند. نوآوری‌های اصلی شامل صورت‌بندی جامع مسئله، بازسازی همزمان مؤلفه‌های کد توربو، مدل‌سازی ابهام ساختاری و تحلیل رابطهٔ نویز با پیچیدگی محاسباتی است. نتایج شبیه‌سازی و تحلیل‌های نظری نشان می‌دهند روش پیشنهادی در شرایط نویزی نیز قادر به بازسازی ساختار مجاز کد است؛ جزئیات عددی پس از تکمیل ارزیابی‌ها در چکیده نهایی خواهد آمد.
}
\keywords{بازسازی کور، کد توربو، کدگذار RSC، درهم‌ریز، مخابرات غیرهمکارانه، شناسایی کور کد کانال}

\graphicalabstract{
\begin{center}
\textit{جای تصویر چکیده گرافیکی: معماری کلی رساله و چارچوب روش پیشنهادی بازسازی کور مشترک.}
\end{center}
}

\similarityresult{
\begin{center}
\textit{جای تصویر نتیجه سامانه تشابه‌یابی (همانندجویی).}
\end{center}
}

\fundingtext{
این رساله تنها با حمایت مالی دانشگاه تهران انجام شده است و حامی مالی دیگری نداشته است.
}

% مشخصات انگلیسی پایان‌نامه
% !TeX root=../main.tex
% مشخصات انگلیسی رساله

%%%%%%%%%%%%%%%%%%%%%%%%%%%%%%%%%%%%
\latinuniversity{University of Tehran}
\latincollege{College of Interdisciplinary Sciences and Technologies}
\latinfaculty{Faculty of Interdisciplinary Sciences and Technologies}
\latindepartment{Department of ??}
\latinsubject{??}
\latinfield{??}
\latintitle{Joint Blind Reconstruction of Turbo Codes with Simultaneous Recovery of RSC Encoders and Inner/Outer Interleavers}
\firstlatinsupervisor{Mehdi Teimouri}
%\secondlatinsupervisor{Second Supervisor}
\firstlatinadvisor{First Advisor}
%\secondlatinadvisor{Second Advisor}
\latinname{Shahryar}
\latinsurname{Toughi}
\latinthesisdate{Month Year}
\latinkeywords{Blind reconstruction, Turbo code, RSC encoder, Interleaver, Non-cooperative communications, Blind channel code identification}
\en-abstract{
Blind identification and reconstruction of channel code structure, particularly turbo codes, is a central problem in non-cooperative communications. This dissertation formulates and addresses the joint reconstruction of RSC generator polynomials together with the inner and outer interleavers, while accounting for trellis termination and channel noise. The proposed approach combines algebraic encoder constraints, statistical constraint testing under noise, and a constrained tree search to recover the unknown parameter set from hard (and, in extensions, soft) observations. The main contributions include a comprehensive problem formulation, simultaneous reconstruction of turbo-code components, modeling of structural ambiguity, and analysis of the relationship between noise, ambiguity, and computational complexity. Simulation results and theoretical analyses indicate that the method can recover admissible code structures under noisy conditions; final quantitative findings will be reported after completing the experimental evaluation.
}


% تنظیمات و تعاریف واژه‌نامه و اختصارات
\input{./tex/glossaries-settings}
% !TeX root=../main.tex
% واژه‌های تخصصی حوزهٔ رساله
\newword{blind-recon}{Blind Reconstruction}{بازسازی کور}{Blind Reconstructions}
\newword{blind-id}{Blind Identification}{شناسایی کور}{Blind Identifications}
\newword{turbo-code}{Turbo Code}{کد توربو}{Turbo Codes}
\newword{rsc-enc}{Recursive Systematic Convolutional Encoder}{کدگذار کانولوشنی سیستماتیک بازگشتی}{Recursive Systematic Convolutional Encoders}
\newword{interleaver}{Interleaver}{درهم‌ریز}{Interleavers}
\newword{trellis}{Trellis}{ترلیس}{Trellises}
\newword{termination}{Trellis Termination}{خاتمهٔ ترلیس}{Trellis Terminations}
\newword{puncturing}{Puncturing}{پانچینگ}{Puncturings}
\newword{awgn-ch}{Additive White Gaussian Noise Channel}{کانال نویز سفید گاوسی جمع‌شونده}{Additive White Gaussian Noise Channels}
\newword{bsc-ch}{Binary Symmetric Channel}{کانال دودویی متقارن}{Binary Symmetric Channels}
\newword{noncoop}{Non-cooperative Communications}{مخابرات غیرهمکارانه}{Non-cooperative Communications}
\newword{ambiguity}{Structural Ambiguity}{ابهام ساختاری}{Structural Ambiguities}
\newword{pruning}{Pruning}{هرس}{Prunings}
\newword{parity}{Parity Constraint}{قید پاریتی}{Parity Constraints}

% !TeX root=../main.tex
% فهرست اختصارات تخصصی
\newacronym{rsc}{RSC}{Recursive Systematic Convolutional}
\newacronym{awgn}{AWGN}{Additive White Gaussian Noise}
\newacronym{bsc}{BSC}{Binary Symmetric Channel}
\newacronym{bpsk}{BPSK}{Binary Phase-Shift Keying}
\newacronym{ldpc}{LDPC}{Low-Density Parity-Check}
\newacronym{bcjr}{BCJR}{Bahl--Cocke--Jelinek--Raviv}
\newacronym{snr}{SNR}{Signal-to-Noise Ratio}
\newacronym{llr}{LLR}{Log-Likelihood Ratio}
\newacronym{ml}{ML}{Maximum Likelihood}
\newacronym{em}{EM}{Expectation--Maximization}


\begin{document}

\pagenumbering{adadi} % یک، دو، ...
% ابتدای درج صفحات مختلف
\coverPage
% بررسی حالت پیش‌نویس
\ifoptiondraft{}{% 
    \besmPage
    \titlePage
} % end ifoptiondraft
\abstractPage
\graphicalAbstractPage
\ifoptiondraft{}{%
    % چنانچه مایل به چاپ صفحات «تقدیم» و «سپاس‌گزاری» نیستید، خطوط زیر را غیرفعال کنید.
    \taghdimPage
    \ghadrdaniPage
    \davaranPage
    \esalatPage
    \similarityPage
    \oathPage
    \fundingPage
    \mojavezPage
} % end ifoptiondraft
% شروع درج فهرست‌ها
\newpage\cleardoublepage
\pagenumbering{harfi} % آ، ب، ...
\tableofcontents \clearpage
% بررسی حالت پیش‌نویس برای بقیه فهرست‌ها
\ifoptiondraft{
    \listoftodos
}{%
    \listoffigures \clearpage
    \listoftables  \clearpage
    \printacronyms
} % end ifoptiondraft


\pagestyle{fancy}
\pagenumbering{arabic} % 1, 2, ...

% !TeX root=../main.tex
\chapter{کلیات تحقیق}
\label{chap:introduction}

\section{مقدمه}
هدف این فصل معرفی مسئله پژوهش، اهمیت و کاربردهای موضوع، و جایگاه رساله در حوزهٔ شناسایی و بازسازی کور%
\footnote{بازسازی کور یعنی بازیابی ساختار ناشناختهٔ کد کانال صرفاً از روی مشاهدات دریافتی، بدون دانش قبلی از پارامترهای فرستنده.}
کدهای کانال است%
\LTRfootnote{Blind reconstruction recovers an unknown channel-code structure from received observations without prior knowledge of the transmitter parameters.}%
. در این فصل از ورود به جزئیات ریاضی الگوریتم خودداری می‌شود و تمرکز بر تبیین روشن مسئله و انگیزهٔ حل آن قرار دارد.

\section{بیان مسئله}
مسئلهٔ محوری این رساله، بازسازی کور مشترک ساختار کد توربو شامل پلی‌نوم‌های کدگذارهای RSC، درهم‌ریز داخلی و درهم‌ریز خارجی، با درنظرگرفتن خاتمهٔ ترلیس و نویز کانال در مخابرات غیرهمکارانه است. برخلاف تشخیص صرف نوع کد، هدف بازیابی پارامترهای ساختاری ناشناخته از روی مشاهدات دریافتی است.
شکل~\ref{fig:turbo-structure-sample} نمای کلی از ساختار موازی یک کدگذار توربو را نشان می‌دهد که مؤلفه‌های اصلی مسئلهٔ بازسازی در این رساله را مشخص می‌کند.

\begin{figure}[ht]
	\centering
	% برای افزودن شکل جدید: فایل تصویر را در پوشه img قرار دهید و نام آن را در دستور زیر عوض کنید.
	\includegraphics[width=0.85\textwidth]{turbo-structure-sample}
	\caption{نمای کلی ساختار موازی کدگذار توربو (نمونه برای الگوبرداری در افزودن شکل‌های بعدی)}
	\label{fig:turbo-structure-sample}
\end{figure}

\section{اهمیت و ضرورت تحقیق}
در سناریوهای غیرهمکارانه، گیرنده معمولاً به مشخصات کامل کدگذاری دسترسی ندارد. بازسازی ساختار کد توربو امکان رمزگشایی، تحلیل ترافیک و توسعهٔ سامانه‌های شنود هوشمند را فراهم می‌کند و از نظر علمی نیز به غنی‌سازی نظریهٔ شناسایی کور کد کانال می‌انجامد.

\section{کاربردهای شناسایی کور کدهای کانال}
کاربردهای عملی شامل مخابرات نظامی و اطلاعاتی، پایش طیف، مهندسی معکوس سامانه‌های مخابراتی، و بازیابی اطلاعات در شرایط فقدان مشخصات فرستنده است. کد توربو به‌عنوان یکی از کدهای رایج استاندارد، مورد مطالعاتی مرکزی این رساله است.

\section{چالش‌های شناسایی کور کدهای توربو}
فضای جست‌وجوی درهم‌ریزها به‌صورت نمایی بزرگ است، نویز قیود جبری را مخدوش می‌کند، و ابهام ساختاری ممکن است منجر به چندین بازسازی مجاز شود. علاوه بر این، خاتمهٔ ترلیس، پانچینگ و همزمان‌سازی فریم پیچیدگی مسئله را افزایش می‌دهند.

\section{اهداف تحقیق}
هدف کلی، ارائهٔ چارچوبی برای شناسایی و بازسازی کور ساختار کد توربو است. اهداف جزئی شامل صورت‌بندی ریاضی پارامترهای مجهول، طراحی الگوریتم بازسازی مشترک، تحلیل پیچیدگی و صحت، و ارزیابی تجربی عملکرد تحت نویز است.

\section{پرسش‌های اصلی تحقیق}
\begin{enumerate}
	\item چگونه می‌توان پلی‌نوم‌های RSC و درهم‌ریزهای داخلی و خارجی را به‌طور مشترک از مشاهدات نویزی بازسازی کرد؟
	\item نقش خاتمهٔ ترلیس و قیود جبری در کاهش فضای جست‌وجو چیست؟
	\item رابطهٔ سطح نویز با ابهام ساختاری و پیچیدگی محاسباتی چگونه است؟
\end{enumerate}

\section{فرضیات و محدودیت‌های تحقیق}
فرضیات پایه شامل مدل کانال AWGN با مدولاسیون BPSK، امکان تبدیل مشاهدهٔ سخت به مدل BSC معادل، و تمرکز بر ساختار موازی کد توربو با خاتمهٔ ترلیس است. محدودیت‌هایی مانند پانچینگ ناشناخته یا عدم همزمانی کامل در فصول توسعه بررسی می‌شوند.

\section{نوآوری‌های تحقیق}
\begin{enumerate}
	\item صورت‌بندی جامع مسئلهٔ شناسایی و بازسازی کور کدهای توربو؛
	\item بازسازی کور پلی‌نوم‌های RSC بدون دانش قبلی؛
	\item بازسازی همزمان درهم‌ریز داخلی و خارجی؛
	\item لحاظ‌کردن خاتمهٔ ترلیس در بازسازی؛
	\item طراحی جست‌وجوی درختی مقید؛
	\item طراحی معیار آماری برای هرس تحت نویز؛
	\item مدل‌سازی و تحلیل ابهام ساختاری؛
	\item تحلیل رابطهٔ نویز، ابهام و پیچیدگی محاسباتی.
\end{enumerate}

\section{روش تحقیق}
روش پژوهش ترکیبی از مدل‌سازی ریاضی، استخراج قیود جبری، آزمون آماری قیود، طراحی الگوریتم جست‌وجوی درختی و ارزیابی با شبیه‌سازی مونت‌کارلو است.

\section{ساختار رساله}
فصل دوم مبانی نظری و پیشینه را مرور می‌کند. فصل سوم مدل سیستم و صورت‌بندی مسئله را ارائه می‌دهد. فصل چهارم روش پیشنهادی را شرح می‌دهد. فصل پنجم به تحلیل نظری و پیچیدگی می‌پردازد. فصل ششم نتایج شبیه‌سازی را گزارش می‌کند. فصل هفتم توسعه‌ها و مطالعات تکمیلی را بررسی می‌کند و فصل هشتم جمع‌بندی و پیشنهادهای آینده را ارائه می‌دهد.
		% فصل اول: کلیات تحقیق
% !TeX root=../main.tex
\chapter{مبانی نظری و پیشینه پژوهش}
\label{chap:background}

\section{مقدمه}
در این فصل ابتدا مبانی علمی لازم برای درک کدهای تصحیح خطا و به‌ویژه کدهای توربو مرور می‌شود؛ سپس مفهوم شناسایی کور و پژوهش‌های مرتبط در کدهای کانولوشنی و توربو دسته‌بندی می‌گردد و در پایان شکاف پژوهشی مشخص می‌شود.

\section{کدهای تصحیح خطا}
\subsection{مفهوم کدگذاری کانال}
کدگذاری کانال با افزودن افزونگی کنترل‌شده، امکان تشخیص و تصحیح خطا در گیرنده را فراهم می‌کند.

\paragraph{کدهای بلوکی}
در کدهای بلوکی، هر بلوک اطلاعات مستقل رمزگذاری می‌شود و ساختار جبری مشخصی دارد.

\paragraph{کدهای کانولوشنی}
کدهای کانولوشنی با حافظهٔ محدود و نمایش ترلیس، پایهٔ بسیاری از کدگذارهای عملی از جمله RSC هستند.

\paragraph{کدهای LDPC}
کدهای بررسی توازن کم‌چگال با گراف دوقسمتی و رمزگشایی تکراری شناخته می‌شوند.

\paragraph{کدهای توربو}
کدهای توربو با ترکیب موازی یا سریال کدگذارهای کانولوشنی و درهم‌ریز، عملکرد نزدیک به ظرفیت شانون را ممکن می‌سازند.

\section{کدهای توربو}
\subsection{تاریخچه و تکامل کدهای توربو}
پیدایش کدهای توربو نقطهٔ عطفی در نظریهٔ کدگذاری بود
\cite{Berrou93}
و به استانداردهای مخابراتی متعددی راه یافت.

\subsection{ساختار Turbo Code}
ساختار متداول شامل دو کدگذار RSC، درهم‌ریز داخلی، و در برخی سامانه‌ها درهم‌ریز خارجی و پانچینگ است.

\subsection{کدگذارهای RSC}
کدگذارهای سیستماتیک بازگشتی پایهٔ بسیاری از طرح‌های توربو هستند و رفتار ترلیس و قیود پاریتی آن‌ها در بازسازی کور نقش کلیدی دارد.

\subsection{پلی‌نوم‌های مولد}
پلی‌نوم‌های بازخورد و پیش‌خور ساختار کدگذار را تعیین می‌کنند و یکی از اهداف اصلی بازسازی کور هستند.

\subsection{حافظه و حالت‌های کدگذار}
حافظهٔ کدگذار تعداد بیت‌های رجیستر و تعداد حالت‌های ترلیس را مشخص می‌کند.

\subsection{ترلیس و گذارهای حالت}
نمایش ترلیس گذارهای مجاز حالت را توصیف می‌کند و مبنای رمزگشایی و استخراج قیود است.

\subsection{خاتمه ترلیس}
بیت‌های خاتمه، ترلیس را به حالت مشخصی می‌برند و طول فریم خروجی را تحت تأثیر قرار می‌دهند؛ در مدل این رساله رابطهٔ طول با خاتمه لحاظ می‌شود.

\subsection{درهم‌ریز داخلی}
درهم‌ریز داخلی ترتیب بیت‌های ورودی کدگذار دوم را تغییر می‌دهد و بازیابی آن از مسائل دشوار بازسازی است.

\subsection{درهم‌ریز خارجی}
درهم‌ریز خارجی بر چینش سریالی خروجی اثر می‌گذارد و باید همراه با سایر پارامترها بازسازی شود.

\subsection{پانچینگ و نرخ کد}
پانچینگ با حذف گزینشی بیت‌های پاریتی نرخ کد را افزایش می‌دهد و شناسایی آن را دشوارتر می‌کند.

\subsection{ساختار سریالی خروجی کد توربو}
چینش بیت‌های سیستماتیک و پاریتی در جریان خروجی، مدل مشاهدهٔ گیرنده را شکل می‌دهد.

\section{مدل کانال و مدل مشاهده}
\subsection{کانال AWGN}
مدل نویز سفید گاوسی جمع‌شونده فرض پایهٔ بسیاری از ارزیابی‌ها است.

\subsection{BPSK}
مدولاسیون دودویی فاز، نگاشت بیت به نماد را ساده می‌کند و با مدل‌های تحلیلی سازگار است.

\subsection{آشکارسازی سخت}
در آشکارسازی سخت، هر نماد به یک بیت قطعی نگاشت می‌شود و مشاهده به دنبالهٔ باینری تبدیل می‌گردد.

\subsection{آشکارسازی نرم}
اطلاعات نرم (مانند LLR) می‌تواند دقت بازسازی را بهبود دهد و در فصل توسعه بررسی می‌شود.

\subsection{مدل کانال معادل BSC}
پس از آشکارسازی سخت، کانال معادل BSC با احتمال خطای بیت مشخص، مدل مشاهدهٔ عملیاتی رساله را تشکیل می‌دهد.

\section{مفهوم شناسایی کور}
در ادبیات، تمایز میان طبقه‌بندی کور، تخمین کور و \gls{blind-recon} اهمیت دارد. طبقه‌بندی نوع کد را مشخص می‌کند، تخمین برخی پارامترها را به‌دست می‌آورد، و بازسازی ساختار کامل (یا مجموعهٔ ساختارهای مجاز) را هدف قرار می‌دهد؛ تمرکز این رساله بر بازسازی است. کدهای \gls{turbo-code} و کدگذارهای \lr{RSC} نمونهٔ مرکزی این مطالعه هستند.

\section{شناسایی کور کدهای کانولوشنی}
روش‌های مبتنی بر رتبه، قیود دوگان، و الگوریتم‌های تکراری برای بازیابی پارامترهای کد کانولوشنی مرور می‌شوند تا زمینه‌ای برای تعمیم به توربو فراهم گردد.

\section{شناسایی کور کدهای توربو}
\subsection{شناسایی پارامترهای کدگذار}
\subsection{بازیابی پلی‌نوم‌ها}
\subsection{بازیابی درهم‌ریز}
\subsection{تشخیص نرخ کد}
\subsection{تشخیص طول فریم}
\subsection{تشخیص حافظه}
\subsection{تشخیص پانچینگ}
\subsection{همزمان‌سازی و تعیین مرز فریم}
در این زیربخش‌ها پژوهش‌های مرتبط با هر پارامتر دسته‌بندی و نقاط قوت و ضعف آن‌ها بیان می‌شود.

\section{روش‌های موجود در ادبیات}
به‌جای فهرست خطی مقالات، روش‌ها در دسته‌هایی مانند شناسایی کدگذار، بازیابی درهم‌ریز، dual-code، BCJR، rank-based، ML/EM و joint reconstruction طبقه‌بندی می‌شوند. جدول~\ref{tab:lit-categories} نمای کلی این دسته‌بندی را نشان می‌دهد.

\begin{table}[ht]
	\caption{دسته‌بندی روش‌های ادبیات در شناسایی/بازسازی کور کد توربو}
	\label{tab:lit-categories}
	\centering
	\begin{tabularx}{0.95\textwidth}{|X|X|}
		\hline
		\textbf{دسته} & \textbf{پارامتر بازیابی‌شدهٔ متداول} \\
		\hline
		شناسایی کدگذار & پلی‌نوم‌ها، حافظه \\
		\hline
		بازیابی درهم‌ریز & جایگزشت \\
		\hline
		dual-code & قیود پاریتی \\
		\hline
		مبتنی بر BCJR & ساختار و درهم‌ریز \\
		\hline
		rank-based & طول، بعد، ساختار \\
		\hline
		ML/EM & تخمین پارامتر \\
		\hline
		بازسازی مشترک & چند مؤلفهٔ همزمان \\
		\hline
	\end{tabularx}
\end{table}

\section{شکاف پژوهشی}
با وجود پیشرفت‌های قابل توجه، بازسازی همزمان ساختار کامل یک کد توربو شامل پارامترهای کدگذار، درهم‌ریز داخلی و درهم‌ریز خارجی، در شرایط مشاهدهٔ ناقص و نویزی، همچنان مسئله‌ای دشوار و تا حد زیادی حل‌نشده است. این رساله همین خلأ را هدف قرار می‌دهد.
		% فصل دوم: مبانی نظری و پیشینه پژوهش
% !TeX root=../main.tex
\chapter{مدل سیستم و صورت‌بندی ریاضی مسئله}
\label{chap:system-model}

\section{مقدمه}
فصل قبل مبانی و پیشینه را مرور کرد؛ در این فصل مسئلهٔ تحقیق به‌صورت دقیق مدل‌سازی و پارامترهای مجهول تعریف می‌شوند.

\section{ساختار عمومی کد توربو}
ساختار موازی متداول شامل دو کدگذار RSC، درهم‌ریز داخلی، و در صورت وجود درهم‌ریز خارجی و الگوی پانچینگ توصیف می‌شود.

\section{مدل کدگذار RSC}
\subsection{مدل چندجمله‌ای}
\subsection{تابع انتقال}
\subsection{نمایش باینری و اکتال}
\subsection{معادلات بازگشتی}
روابط بازگشتی میان بیت سیستماتیک و پاریتی، قیود جبری لازم برای بازسازی را فراهم می‌کنند.

\section{مدل درهم‌ریز داخلی}
درهم‌ریز داخلی به‌صورت جایگزشت روی بلوک اطلاعات مدل می‌شود و فضای حالت آن در حالت کلی بسیار بزرگ است.

\section{مدل درهم‌ریز خارجی}
درهم‌ریز خارجی چینش بیت‌های خروجی را تعیین می‌کند و باید در مجموعهٔ پارامترهای مجهول لحاظ شود.

\section{مدل پانچینگ}
الگوی پانچینگ نرخ مؤثر کد را تغییر می‌دهد؛ در مدل پایه ممکن است معلوم فرض شود و در فصل توسعه به‌صورت ناشناخته بررسی گردد.

\section{مدل خاتمه ترلیس}
با درنظرگرفتن بیت‌های خاتمه، طول فریم پیش از درهم‌ریز خارجی می‌تواند از رابطه‌ای مانند
\begin{equation}
	N = 3L + 4m
	\label{eq:frame-length}
\end{equation}
به‌دست آید که در آن \(L\) طول بلوک اطلاعات و \(m\) حافظهٔ کدگذار است.

\section{مدل کانال}
\subsection{AWGN}
\subsection{BPSK}
\subsection{BSC معادل}
پس از آشکارسازی سخت، مشاهده به‌صورت دنباله‌های باینری نویزی در مدل BSC توصیف می‌شود.

\section{مدل دریافت}
فرض‌های مشاهده شامل مشخص‌بودن یا نبودن مرز فریم، نوع آشکارسازی سخت/نرم، تعداد فریم‌های دریافتی \(F\)، نرخ خطای کانال، و ناشناخته‌بودن ساختار کد است.

\section{تعریف پارامترهای مجهول}
مجموعهٔ پارامترهای مجهول به‌صورت
\begin{equation}
	\theta = \{m, g_0(D), g_1(D), \pi_i, \pi_o, L\}
	\label{eq:theta}
\end{equation}
تعریف می‌شود که شامل حافظه، پلی‌نوم‌های مولد، درهم‌ریزهای داخلی و خارجی، و طول بلوک است.

\section{تعریف مسئله بازسازی کور}
با دریافت مشاهدات \(R\)، هدف تخمین مجموعهٔ مقادیر مجاز \(\theta\) است. خروجی الگوریتم لزوماً یک پاسخ یکتا نیست و ممکن است مجموعه‌ای از بازسازی‌های ساختاریاً معادل (ابهام ساختاری) به‌دست آید.
		% فصل سوم: مدل سیستم و صورت‌بندی ریاضی مسئله
% !TeX root=../main.tex
\chapter{روش پیشنهادی بازسازی کور مشترک}
\label{chap:proposed-method}

\section{مقدمه}
این فصل قلب رساله است و معماری الگوریتم بازسازی کور مشترک کد توربو را از ایدهٔ اصلی تا تحلیل درستی شرح می‌دهد.

\section{ایده اصلی روش پیشنهادی}
ایدهٔ محوری، ترکیب قیود جبری RSC با آزمون آماری تحت نویز و جست‌وجوی درختی مقید است تا فضای جایگزشت‌ها و پلی‌نوم‌ها به‌صورت تدریجی هرس شود.

\section{تعریف فضای جست‌وجو}
\subsection{فضای پلی‌نوم‌ها}
\subsection{فضای حافظه}
\subsection{فضای درهم‌ریزها}
\subsection{محدودیت‌های طول فریم}
فضای خام درهم‌ریز از مرتبهٔ \(N!\) است و جست‌وجوی کامل عملاً ناممکن است؛ بنابراین کاهش فضا با قیود ضروری است.

\section{استخراج قیود جبری RSC}
\subsection{رابطه recursive}
\subsection{قیود parity}
\subsection{قیود polynomial realizability}
این قیود نامزدهای ناسازگار را در مراحل اولیه حذف می‌کنند.

\section{آزمون آماری قیود}
\subsection{احتمال برقراری constraint}
\subsection{احتمال نقض constraint}
\subsection{مدل Binomial}
\subsection{تعیین آستانه رد}
\subsection{کنترل احتمال پذیرش کاذب}
تحت نویز، برقراری دقیق قیود تضمین نمی‌شود؛ بنابراین آزمون آماری مبتنی بر مدل دوجمله‌ای برای پذیرش یا رد فرضیه‌ها به کار می‌رود.

\section{معماری جست‌وجوی درختی}
\subsection{تعریف Node}
\subsection{تعریف Depth}
\subsection{تعریف Hypothesis}
\subsection{Priority Queue}
\subsection{Best-first Search}
جست‌وجوی بهترین‌اول با صف اولویت، بسط فرضیه‌های امیدوارکننده‌تر را در اولویت قرار می‌دهد.

\section{فاز اول بازسازی}
\subsection{بازسازی \(U_1\)}
\subsection{بازسازی \(P_1\)}
\subsection{بازسازی \(q(D)\)}
\subsection{استخراج \(g_0(D)\)}
\subsection{استخراج \(g_1(D)\)}
\subsection{تعیین حافظه \(m\)}
\subsection{کنترل طول فریم}
فاز اول عمدتاً مؤلفه‌های مرتبط با کدگذار اول و طول سازگار را بازیابی می‌کند.

\section{فاز دوم بازسازی}
\subsection{بازسازی \(U_2\)}
\subsection{بازسازی \(P_2\)}
\subsection{قیود interleaver}
\subsection{بازیابی \(\pi_i\)}
\subsection{بازیابی \(\pi_o\)}
فاز دوم با اتکا به نتایج فاز اول، دنباله‌ها و درهم‌ریزهای داخلی و خارجی را بازسازی می‌کند.

\section{راهبردهای Pruning}
\subsection{Parity-consistency}
\subsection{Polynomial realizability}
\subsection{Interleaver consistency}
\subsection{Encoder-domain constraint}
\subsection{Heuristic pruning}
هرس چندلایه مانع انفجار شاخه‌ها می‌شود و باید احتمال حذف پاسخ واقعی را کنترل کند.

\section{الگوریتم نهایی}
در این بخش شبه‌کد کامل الگوریتم دوفازی همراه با ورودی/خروجی و معیار توقف ارائه می‌شود.

\begin{algorithm}[ht]
\caption{اسکلت الگوریتم بازسازی کور مشترک کد توربو}
\label{alg:joint-blind}
\begin{algorithmic}[1]
\REQUIRE مشاهدات \(R\)، محدودهٔ جست‌وجوی \(m\) و طول، آستانه‌های آماری
\ENSURE مجموعهٔ پاسخ‌های مجاز \(\widehat{\Theta}\)
\STATE مقداردهی صف اولویت با فرضیه‌های اولیهٔ فاز اول
\WHILE{صف خالی نیست}
	\STATE انتخاب بهترین فرضیه
	\STATE آزمون قیود جبری و آماری
	\IF{فرضیه رد شود}
		\STATE هرس و ادامه
	\ENDIF
	\STATE بسط فرزندان و به‌روزرسانی صف
\ENDWHILE
\STATE اجرای فاز دوم برای نامزدهای باقی‌مانده و تشکیل \(\widehat{\Theta}\)
\RETURN \(\widehat{\Theta}\)
\end{algorithmic}
\end{algorithm}

\section{تحلیل درستی الگوریتم}
\subsection{شرط لازم}
\subsection{شرط کافی}
\subsection{احتمال حذف جواب واقعی}
\subsection{احتمال پذیرش جواب کاذب}
\subsection{کامل بودن یا عدم کامل بودن جست‌وجو}
در این بخش، خواص نظری الگوریتم نسبت به مقالهٔ پایه گسترش می‌یابد تا برای سطح رسالهٔ دکتری کفایت استدلالی داشته باشد.
		% فصل چهارم: روش پیشنهادی بازسازی کور مشترک
% !TeX root=../main.tex
\chapter{تحلیل نظری و پیچیدگی محاسباتی}
\label{chap:complexity}

\section{مقدمه}
این فصل پیچیدگی فضای جست‌وجو، اثر قیود و نویز بر هزینهٔ محاسباتی، و کران‌های عملی روش پیشنهادی را به‌صورت مستقل از ارزیابی شبیه‌سازی تحلیل می‌کند.

\section{پیچیدگی فضای جست‌وجوی کامل}
ابعاد فضای کامل شامل ترکیب فضای پلی‌نوم، حافظه و جایگزشت‌ها توصیف می‌شود.

\section{پیچیدگی حالت brute-force}
جست‌وجوی کامل درهم‌ریز از مرتبهٔ فاکتوریل است و عملاً غیرقابل اجراست؛ این بخش کران پایین محاسباتی را روشن می‌کند.

\section{پیچیدگی جست‌وجوی محدودشده}
با اعمال قیود جبری و آماری، فضای مؤثر به‌صورت قابل توجهی کاهش می‌یابد؛ مدل پیچیدگی الگوریتم مقید در اینجا ارائه می‌شود.

\section{تحلیل Branching Factor}
ضریب انشعاب درخت جست‌وجو در اعماق مختلف و وابستگی آن به نوع قیود بررسی می‌شود.

\section{اثر قیود بر کاهش فضای جست‌وجو}
هر قید (پاریتی، تحقق‌پذیری پلی‌نوم، سازگاری درهم‌ریز) سهم جداگانه‌ای در هرس دارد که باید کمی‌سازی شود.

\section{رابطه پیچیدگی با نویز}
با افزایش نویز، تعداد فرضیه‌هایی که از هرس عبور می‌کنند افزایش می‌یابد و رفتار جست‌وجو به سمت رشد نمایی میل می‌کند؛ این پدیده از یافته‌های مهم تحلیل است.

\section{اثر حافظه کدگذار}
افزایش \(m\) هم فضای پلی‌نوم و هم طول خاتمه را بزرگ‌تر می‌کند و هزینه را بالا می‌برد.

\section{اثر طول interleaver}
طول بزرگ‌تر جایگزشت، فضای نامزدها و زمان بسط گره‌ها را افزایش می‌دهد.

\section{تحلیل زمان اجرا}
مدل زمان اجرا برحسب تعداد گره‌های بسط‌یافته، آزمون‌های آماری و عملیات هرس ارائه می‌شود.

\section{کران‌های عملی روش}
شرایطی که در آن روش از نظر زمان و حافظه عملی باقی می‌ماند، همراه با رژیم‌های ناکارآمد، جمع‌بندی می‌گردد.
		% فصل پنجم: تحلیل نظری و پیچیدگی محاسباتی
% !TeX root=../main.tex
\chapter{نتایج شبیه‌سازی و ارزیابی عملکرد}
\label{chap:results}

\section{محیط شبیه‌سازی}
بسته‌های نرم‌افزاری، سخت‌افزار اجرا، و چارچوب تولید داده و نویز در این بخش توصیف می‌شوند.

\section{پارامترهای شبیه‌سازی}
پارامترهای پایه شامل حافظهٔ \(m\)، طول \(L\)، تعداد فریم \(F\)، احتمال خطای بیت \(P_b\)، نوع درهم‌ریز، نوع کد، خاتمه، پانچینگ و نوع مشاهدهٔ سخت/نرم است. نمونه‌های موردی مانند کدهای \((13,15)\)، \((15,17)\) و \((133,171)\) و طول‌های \(L\in\{50,100,300,500\}\) در ارزیابی لحاظ می‌شوند.

\section{معیارهای ارزیابی}
\subsection{Reconstruction Success Rate}
\subsection{False Reconstruction Rate}
\subsection{Number of admissible solutions}
\subsection{Structural distance}
\subsection{Runtime}
\subsection{Number of expanded nodes}
\subsection{Number of pruned nodes}
این معیارها فراتر از دقت صرف، جنبهٔ ابهام و هزینهٔ محاسباتی را نیز پوشش می‌دهند.

\section{اثر SNR / \(P_b\)}
منحنی‌های موفقیت و پیچیدگی برحسب سطح نویز گزارش می‌شوند.

\section{اثر طول interleaver}
تأثیر طول جایگزشت بر نرخ موفقیت و زمان اجرا بررسی می‌شود.

\section{اثر حافظه کدگذار}
عملکرد برای مقادیر مختلف \(m\) مقایسه می‌گردد.

\section{اثر تعداد فریم‌های دریافتی}
افزایش \(F\) معمولاً آمار قیود را پایدارتر می‌کند؛ این بخش حساسیت را می‌سنجد.

\section{اثر termination}
نقش مدل خاتمه در صحت طول فریم و بازسازی ساختار تحلیل می‌شود.

\section{اثر نوع interleaver}
درهم‌ریزهای تصادفی و ساخت‌یافته از نظر دشواری بازیابی مقایسه می‌شوند.

\section{اثر hard و soft observation}
تفاوت مشاهدهٔ سخت و نرم در دقت و هزینه بررسی می‌شود.

\section{تحلیل ambiguity}
حتی در حالت بدون نویز ممکن است چند بازسازی معادل وجود داشته باشد؛ این بخش توزیع تعداد جواب‌های مجاز را تحلیل می‌کند.

\section{مقایسه با روش‌های موجود}
مقایسه در چهار محور انجام می‌شود: گسترهٔ بازسازی، دانش پیشین موردنیاز، مقاومت به نویز، و پیچیدگی محاسباتی.

\section{تحلیل آماری Monte Carlo}
تنظیمات تکرار مونت‌کارلو، فاصلهٔ اطمینان و آزمون‌های آماری مرتبط گزارش می‌شوند.

\section{بحث نتایج}
یافته‌های کلیدی، محدودیت‌های مشاهده‌شده و دلالت‌های آن‌ها برای طراحی الگوریتم جمع‌بندی می‌گردند.
		% فصل ششم: نتایج شبیه‌سازی و ارزیابی عملکرد
% !TeX root=../main.tex
\chapter{توسعه روش پیشنهادی و مطالعات تکمیلی}
\label{chap:extensions}

\section{مقدمه}
این فصل رساله را از یک مقالهٔ گسترش‌یافته فراتر می‌برد و تعمیم‌ها و مطالعات تکمیلی حول هستهٔ بازسازی کور مشترک را معرفی می‌کند.

\section{بازسازی در شرایط عدم همزمانی}
حالتی بررسی می‌شود که مرز دقیق نماد یا فریم در دسترس نباشد و همزمان‌سازی جزئی از مسئلهٔ بازسازی گردد.

\section{تخمین طول فریم ناشناخته}
وقتی \(L\) یا \(N\) از پیش معلوم نیست، راهبردهای تخمین طول و اتصال آن به جست‌وجوی ساختاری مطرح می‌شود.

\section{تخمین نرخ کد}
شناسایی نرخ مؤثر و ارتباط آن با الگوی افزونگی/پانچینگ بررسی می‌گردد.

\section{بازسازی تحت puncturing ناشناخته}
الگوی پانچینگ به مجموعهٔ پارامترهای مجهول افزوده می‌شود و اثر آن بر قیود و پیچیدگی تحلیل می‌شود.

\section{بازسازی با soft information}
استفاده از اطلاعات نرم برای بهبود آزمون قیود و کاهش پذیرش کاذب مطالعه می‌شود.

\section{بازسازی با تعداد فریم محدود}
عملکرد روش در رژیم دادهٔ کم و راهبردهای پایدارسازی آمار قیود بررسی می‌گردد.

\section{بازسازی در شرایط کانال متغیر}
کانال ناایستا یا تغییر آهستهٔ SNR چالش‌هایی برای آستانه‌های ثابت ایجاد می‌کند که در اینجا تحلیل می‌شود.

\section{استفاده از روش‌های هوشمند برای هدایت جست‌وجو}
رهیافت‌های مکاشفه‌ای برای اولویت‌بندی فرضیه‌ها و کاهش بسط غیرضروری معرفی می‌شوند.

\section{مقایسه Best-first، Beam Search و روش‌های دیگر}
راهبردهای جایگزین پیمایش درخت از نظر دقت، حافظه و زمان مقایسه می‌گردند.

\section{امکان استفاده از یادگیری ماشین برای pruning}
امکان جایگزینی یا تکمیل قواعد هرس کلاسیک با مدل‌های یادگیری برای پیش‌بینی نامزدهای امیدوارکننده بررسی می‌شود.
		% فصل هفتم: توسعه روش و مطالعات تکمیلی
% !TeX root=../main.tex
\chapter{جمع‌بندی و پیشنهادهای آینده}
\label{chap:conclusion}

\section{جمع‌بندی تحقیق}
این رساله مسئلهٔ شناسایی و بازسازی کور ساختار کدهای توربو در مخابرات غیرهمکارانه را صورت‌بندی کرد و روشی برای بازسازی مشترک پلی‌نوم‌های RSC و درهم‌ریزهای داخلی و خارجی با لحاظ خاتمهٔ ترلیس و نویز ارائه داد.

\section{دستاوردهای اصلی}
دستاوردها شامل چارچوب ریاضی پارامترهای مجهول، الگوریتم جست‌وجوی درختی مقید دوفازی، معیارهای آماری هرس، تحلیل ابهام ساختاری، و ارزیابی نظری/تجربی پیچیدگی و عملکرد است.

\section{پاسخ به پرسش‌های تحقیق}
در این بخش، پاسخ‌های مستند به پرسش‌های مطرح‌شده در فصل اول بر اساس نتایج فصول چهارم تا ششم جمع‌بندی می‌شود.

\section{محدودیت‌های روش}
محدودیت‌هایی مانند حساسیت به نویز بالا، هزینهٔ محاسباتی در طول‌های بزرگ، وابستگی به فرضیات مدل مشاهده، و ابهام ساختاری باقی‌مانده بیان می‌گردد.

\section{کاربردهای عملی}
دلالت‌های روش برای سامانه‌های غیرهمکارانه، پایش طیف و مهندسی معکوس کد کانال مرور می‌شود.

\section{مسیرهای پژوهشی آینده}
پیشنهادهای آینده شامل تعمیم به کدهای مرتبط، بازسازی برخط، بهره‌گیری قوی‌تر از اطلاعات نرم، یادگیری برای هرس، و ارزیابی روی داده‌های واقعی میدان است.
		% فصل هشتم: جمع‌بندی و پیشنهادهای آینده

% منابع علمی (سبک IEEE)
\pagestyle{empty}
{
\small
\onehalfspacing
\bibliographystyle{ieeetr-fa}
\bibliography{./tex/MyReferences}
}

\pagestyle{fancy}

% واژه‌نامه‌ها پیش از پیوست‌ها (مطابق قالب دانشکده)
\onehalfspacing
\cleardoublepage
\printglossary

\appendix
\clearpage
% ثبت تغییر پیشوند فهرست مطالب پیش از اولین پیوست
\addtocontents{toc}{%
    \protect\renewcommand\protect\cftchappresnum{\appendixname~}%
    \protect\setlength{\cftchapnumwidth}{\mylenapp}}
\null\clearpage

% !TeX root=../main.tex
\chapter{مطالب تکمیلی نظری}
\label{app:theory}

این پیوست برای اثبات‌های تکمیلی، مشتق‌گیری‌های طولانی و جزئیات نظری‌ای در نظر گرفته شده است که برای روانی متن اصلی در فصول سوم تا پنجم گنجانده نشده‌اند.
		% پیوست الف: مطالب تکمیلی
% !TeX root=../main.tex
\chapter{جداول و نتایج تکمیلی}
\label{app:extra-results}

این پیوست برای جداول پارامتر، نمودارهای تکمیلی مونت‌کارلو و نتایجی است که به‌دلیل حجم در فصل ششم نیامده‌اند.
		% پیوست ب: نتایج و جداول تکمیلی
% !TeX root=../main.tex
\chapter{جزئیات الگوریتمی و پیاده‌سازی}
\label{app:algo-details}

این پیوست برای شبه‌کدهای تفصیلی، ساختار داده‌ها، تنظیمات پیاده‌سازی و نکات تکرارپذیری الگوریتم فصل چهارم اختصاص دارد.
		% پیوست پ: جزئیات الگوریتمی

\pagenumbering{tartibi} % اول، دوم، ...

\cleardoublepage
\printindex

\begin{latin}
\baselineskip=.6cm
\latinabstract
\latinTitlePage
\end{latin}

\end{document}
